
\section{Implementation}
% 实现
\label{sec:implementation}

We implemented \sys on Linux 6.15 with kernel extensions for hook registration and policy verification.
% 我们在 Linux 6.15 上实现了 \sys，通过内核扩展支持钩子注册与策略验证。
Instrumentation in NVIDIA's open GPU kernel modules exposes memory and scheduling state transitions to policy programs (\S\ref{sec:impl-host}).
% NVIDIA 开放 GPU 内核模块中的插桩向策略程序暴露内存与调度状态转移（\S\ref{sec:impl-host}）。
The user-space loader and JIT build on bpftime~\cite{bpftime}, with an LLVM backend translating the restricted eBPF instruction subset into PTX (\S\ref{sec:impl-gpu}).
% 用户态加载器与 JIT 基于 bpftime~\cite{bpftime}，由 LLVM 后端将受限的 eBPF 指令子集翻译为 PTX（\S\ref{sec:impl-gpu}）。
Although our implementation targets NVIDIA GPUs, the host-side design aligns with generic Linux abstractions (HMM/\texttt{migrate\_vma}~\cite{rocm-gpu-memory,linux-gpusvm-doc}, DRM scheduler~\cite{drm-gpu-sched,amdgpu-userq-doc}), and the device-side can be extended to generate vendor-neutral IR~\cite{spirv-ir-rfc} for future portability.
% 尽管我们的实现面向 NVIDIA GPU，但主机端设计遵循通用 Linux 抽象（HMM/\texttt{migrate\_vma}~\cite{rocm-gpu-memory,linux-gpusvm-doc}、DRM 调度器~\cite{drm-gpu-sched,amdgpu-userq-doc}），设备端也可扩展为生成厂商中立的 IR~\cite{spirv-ir-rfc}，以支持未来的可移植性。

\subsection{Host-side Integration}
% 主机端集成
\label{sec:impl-host}
\paragraph{Memory State-Transition Hooks.}
% 内存状态变迁钩子
\sys extends the NVIDIA Open GPU OS Kernel Modules~\cite{nvidia-open-gpu} with hooks for memory allocation, page faults, eviction-candidate selection, and proactive migration.
% \sys 在 NVIDIA 开放 GPU OS 内核模块~\cite{nvidia-open-gpu} 中添加内存分配、缺页、驱逐候选选择和主动迁移钩子。
At each hook, the driver constructs a context (e.g., 2MB-aligned region range, fault address, block size) and invokes the attached eBPF handler.
% 在每个钩子处，驱动构造一个上下文（例如 2MB 对齐的区域范围、缺页地址、块大小），并调用挂载的 eBPF 处理函数。
Eviction hooks operate at 2MB granularity; prefetch hooks operate at page granularity and can prefetch as many pages as needed.
% 驱逐钩子以 2MB 粒度工作；预取钩子以页粒度工作，可按需预取任意数量的页。
Handlers request asynchronous migration and preemption through restricted driver interfaces, with deferred work executed on background workqueues.
% 处理函数通过受限的驱动接口请求异步迁移与抢占，延后工作在后台工作队列上执行。
The driver maintains a doubly-linked eviction list; handlers may reorder regions via provided helpers, but the driver retains eviction authority under memory pressure.
% 驱动维护一个双向链接的驱逐列表；处理函数可以通过提供的辅助函数重排区域，但在内存压力下驱逐决定权仍归驱动所有。

\paragraph{Scheduling Hooks.}
% 调度钩子
Scheduling hooks observe the initialization and teardown of Time-Slice Groups (TSGs).
% 调度钩子观测 Time-Slice Group（TSG）的初始化与销毁。
Policies configure TSG priority, timeslice duration, and runlist interleaving through restricted driver interfaces.
% 策略通过受限的驱动接口配置 TSG 优先级、时间片长度和运行列表交错。
Policies triggered from uprobes or tracepoints can also preempt running GPU kernels precisely with preemption-request kfuncs.
% 由 uprobe 或 tracepoint 触发的策略还可以通过抢占请求 kfunc 精确地抢占正在运行的 GPU 内核。

\subsection{GPU Runtime}
% GPU 运行时
\label{sec:impl-gpu}

\paragraph{Verification and JIT Compilation.}
% 验证与 JIT 编译
Kernel-resident host programs use Linux's eBPF verifier; device programs use PREVAIL~\cite{prevail-verifier} with GPU helper and map models, followed by the uniformity analysis and SIMT checks in \S\ref{sec:verification}.
% 驻留内核的主机程序使用 Linux 的 eBPF 验证器；设备程序使用带 GPU helper 与 map 模型的 PREVAIL~\cite{prevail-verifier}，随后执行 \S\ref{sec:verification} 中的一致性分析与 SIMT 检查。
The user-space backend compiles accepted device bytecode into PTX.
% 用户态后端将通过验证的设备字节码编译为 PTX。
During JIT compilation, we inline helper functions and map accesses.
% 在 JIT 编译过程中，我们内联辅助函数和 map 访问。
Unlike the workshop version~\cite{egpu}, which relies on atomic synchronization, \sys{} avoids cross-SM primitives to prevent hardware stalls.
% 与依赖原子同步的研讨会版本~\cite{egpu}不同，gpubpf 避免使用跨 SM 原语，以防止硬件停顿。

\paragraph{Device-Side State-Transition Hooks.}
% 设备端状态变迁钩子
\sys hooks device-side state transitions (GPU kernel launch, memory access, execution-phase boundaries) by using ptrace to pause applications, intercepting CUDA runtime APIs, extracting GPU kernel PTX, inserting binary trampolines, and reloading instrumented GPU kernels without application recompilation or restart.
% gpubpf 通过 ptrace 暂停应用、拦截 CUDA 运行时 API、提取 GPU 内核 PTX、插入二进制 trampoline 并重新加载插桩后的 GPU 内核，对设备端状态变迁（GPU 内核启动、内存访问、执行阶段边界）进行挂钩，且无需重新编译或重启应用。
Trampolines insert prologue/epilogue logic for SIMT register preservation.
% Trampoline 插入用于 SIMT 寄存器保存的前导/收尾逻辑。
Each trampoline aggregates per-lane inputs, selects an active warp leader to execute the policy once, and broadcasts its decisions to the participating lanes.
% 每个 trampoline 聚合逐 lane 输入，选择一个活跃的 warp leader 执行一次策略，并将决策广播给参与的 lane。
Aggregation reduces duplicate work within a warp; it does not make total instrumentation cost independent of block count.
% 聚合减少了 warp 内的重复工作，但并不使插桩总成本独立于线程块数量。
Total cost still depends on the number of executed hooks and their map accesses or output records, as illustrated by the complete-record and GPU-local-buffer variants in \S\ref{sec:rq2-cost}.
% 总成本仍取决于实际执行的钩子数量及其 map 访问或输出记录，\S\ref{sec:rq2-cost} 中的完整记录流与 GPU 本地缓冲变体展示了这一点。
For block-level scheduling, \sys implements policies using persistent GPU kernels that poll and claim logical work units; on Blackwell GPUs, we use cluster launch control APIs~\cite{cuda-clc} to schedule thread blocks.
% 对于块级调度，gpubpf 使用常驻 GPU 内核（轮询并认领逻辑工作单元）来实现策略；在 Blackwell GPU 上，我们使用 cluster launch control API~\cite{cuda-clc} 来调度线程块。
For memory observation, \sys attaches eBPF handlers at memory instructions, GPU kernel function boundaries, and tracepoints to capture access patterns and issue policy hints to the host driver via cross-layer maps (e.g., device L2 prefetch can trigger a host callback for extended prefetch).
% 对于内存观测，gpubpf 在内存指令、GPU 内核函数边界和 tracepoint 处挂载 eBPF 处理函数，以捕获访问模式并通过跨层 map 向主机端驱动下发策略提示（例如，设备端 L2 预取可触发主机端回调以进行扩展预取）。

\paragraph{Binaries without PTX}
% 不含 PTX 的二进制程序。
Host-side policies do not require application PTX; for device-side observation, we additionally implement a bounded SASS-injection prototype using NVBit~\cite{villa2019nvbit}.
% 主机端策略不依赖应用 PTX；对于设备端观测，我们还使用 NVBit~\cite{villa2019nvbit} 实现了一个范围受限的 SASS 注入原型。
The prototype verifies a BPF program, compiles it into a device function, and embeds its assembled SASS in an NVBit tool that inserts calls before the target kernel's EXIT instructions.
% 该原型验证 BPF 程序，将其编译为设备函数，再把汇编后的 SASS 嵌入 NVBit 工具，由工具在目标内核的 EXIT 指令前插入调用。
On an RTX 5090, an existing vector-add application compiled with cubin only executes the BPF callback on all 100,352 launched threads while retaining its expected output.
% 在 RTX 5090 上，仅包含 cubin 的既有向量加法程序在全部 100,352 个启动线程上执行了 BPF 回调，同时保持了预期计算结果。
The callback runs inside the target kernel, not as a separately launched companion kernel.
% 回调在目标内核内部运行，而不是作为单独启动的伴随内核运行。
This result covers one context, one selected kernel and an eight-byte per-thread output; it does not establish general helper/map support, late attachment, or performance parity with the PTX path.
% 该结果覆盖单个上下文、单个选定内核和逐线程八字节输出；它并未证明通用 helper/map 支持、运行中挂载或与 PTX 路径的性能相当。
% Evidence: gpu_ext/workloads/sass-kretprobe/results/sass-exit-575-20260907-01/.

\paragraph{Hierarchical Maps and Cross-Domain State}
% 分层 Map 与跨域状态

We realize cross-layer maps (\S\ref{sec:cross-layer-maps}) by placing shards in GPU-accessible memory tiers: cold state in host DRAM, hot state in GPU global memory, and per-warp data in GPU shared memory.
% 我们通过将分片（shard）放置于 GPU 可访问的不同内存层级来实现跨层 map（见 \S\ref{sec:cross-layer-maps}）：冷状态放在主机 DRAM，热状态放在 GPU 全局内存，逐 warp 数据放在 GPU 共享内存。
A runtime daemon asynchronously flushes GPU-local shards to host-visible canonical instances, providing coherent snapshots without synchronization overhead.
% 运行时守护进程异步地将 GPU 本地分片刷回主机端可见的规范实例，在无需同步开销的情况下提供一致性快照。
Shard updates occur at warp granularity through warp-leader threads without GPU atomic operations: handlers accumulate counters in warp-local registers, aggregate per warp, and store into GPU-local shards for periodic host flush.
% 分片更新以 warp 粒度通过 warp leader 线程完成，且不依赖 GPU 原子操作：处理函数在 warp 本地寄存器中累积计数器，按 warp 聚合后写入 GPU 本地分片，再由主机端定期刷回。
